\begin{apartats}

\apartat{1}
Calcula els límits següents:
\begin{graella}{3}
  \sa \lim_{x\to+\infty}\frac{3x^2+1}{x-4} &
  \sa \lim_{x\to-\infty}\frac{1-x^4}{2x^4-x^2+5} &
  \sa \lim_{x\to+\infty}\frac{x^3+2x}{3^x}
\end{graella}

\begin{solucio}
i) El grau del numerador supera el del denominador i el quocient dels coeficients
principals és positiu: $+\infty$.\\
ii) Mateix grau: quocient dels coeficients principals, $-1/2$.\\
iii) L'exponencial és un infinit d'ordre superior a qualsevol potència: $0$.
\end{solucio}

\apartat{0,75}
Calcula, si existeixen, els límits següents. Si algun no existeix, justifica-ho amb
els límits laterals.
\begin{graella}{2}
  \sa \lim_{x\to-2}\frac{x^2+x-2}{x^2+3x+2} &
  \sa \lim_{x\to1}\frac{x^2-1}{x^2-2x+1}
\end{graella}

\begin{solucio}
i) $\dfrac{(x+2)(x-1)}{(x+1)(x+2)}=\dfrac{x-1}{x+1}\to\dfrac{-3}{-1}=3$.\\
ii) $\dfrac{(x-1)(x+1)}{(x-1)^2}=\dfrac{x+1}{x-1}$. Els laterals valen $-\infty$ i
$+\infty$, per tant el límit \textbf{no existeix}.
\end{solucio}

\apartat{0,75}
Donada la funció
\[
g(x)=\begin{cases} x^2+1 & \si{x<2},\\[4pt] \dfrac{12}{x+1} & \si{x\ge 2},\end{cases}
\]
determina:
\begin{graella}{3}
  \sa \lim_{x\to-1}g(x) & \sa \lim_{x\to2}g(x) & \sa \lim_{x\to+\infty}g(x)
\end{graella}

\begin{solucio}
i) Com que $-1<2$, hi actua la branca $x^2+1$: el límit val $2$. (La branca
$12/(x+1)$ no hi intervé.)\\
ii) Laterals $5$ i $4$: el límit \textbf{no existeix}.\\
iii) $12/(x+1)\to 0$.
\end{solucio}

\end{apartats}
